#import "@local/bootlin:0.1.0": *

#import "@local/bootlin-yocto:0.1.0": *

#import "@local/bootlin-utils:0.1.0": *

#import "../../typst/local/themeBootlin.typ": *

#import "../../typst/local/common.typ": *

#show: bootlin-theme.with( aspect-ratio: "16-9", config-common(handout: "handout" in sys.inputs and sys.inputs.handout == "1", ))

#show raw.where(block: true): set block(fill: luma(240), inset: 1em, radius: 0.5em, width: 100\%)

#show raw.where(block: true): set text(size: 11pt)

#show raw.where(block: false): r => text(fill: color-link)[#r]

#show raw.where(lang: "c", block: true): set block(fill: luma(240), inset: 0.4em, radius: 0.5em, width: 95\%, breakable: true, above: 12pt, below: 12pt)

#show raw.where(lang: "c", block: true): set text(11pt)

#show raw.where(lang: "console", block: true):set block(fill: luma(240), inset: 0.4em, radius: 0.5em, width: 95\%, breakable: true, above: 6pt)

#show raw.where(lang:"console", block: true): set text(12pt)

 == Contents

=== Root filesystem organization
  \begin{itemize}
  \item The organization of a Linux root filesystem in terms of
    directories is well-defined by the {\bf Filesystem Hierarchy
      Standard}
  \item \url{https://refspecs.linuxfoundation.org/fhs.shtml}
  \item Most Linux systems conform to this specification
    \begin{itemize}
    \item Applications expect this organization
    \item It makes it easier for developers and users as the
      filesystem organization is similar in all systems
    \end{itemize}
  \end{itemize}


=== Important directories (1)
  \begin{description}
  \item[/bin] Basic programs
  \item[/boot] Kernel images, configurations and initramfs
    (only when the kernel is loaded from a
    filesystem, not common on non-x86 architectures)
  \item[/dev] Device files (covered later)
  \item[/etc] System-wide configuration
  \item[/home] Directory for the users home directories
  \item[/lib] Basic libraries
  \item[/media] Mount points for removable media
  \item[/mnt] Mount point for a temporarily mounted filesystem
  \item[/proc] Mount point for the proc virtual filesystem
  \end{description}


=== Important directories (2)
  \begin{description}
  \item[/root]Home directory of the `root` user
  \item[/run]Run-time variable data (previously `/var/run`)
  \item[/sbin]Basic system programs
  \item[/sys]Mount point of the sysfs virtual filesystem
  \item[/tmp]Temporary files
  \item[/usr]
    \begin{description}
    \item[/usr/bin]Non-basic programs
    \item[/usr/lib]Non-basic libraries
    \item[/usr/sbin]Non-basic system programs
    \end{description}
  \item[/var] Variable data files, for system services. This includes spool directories and
    files, administrative and logging data, and transient and
    temporary files
  \end{description}


=== Separation of programs and libraries
  \small
  \begin{itemize}
  \item Basic programs are installed in `/bin` and `/sbin`
    and basic libraries in `/lib`
  \item All other programs are installed in `/usr/bin` and
    `/usr/sbin` and all other libraries in `/usr/lib`
  \item In the past, on UNIX systems, `/usr` was very often
    mounted over the network, through NFS
  \item In order to allow the system to boot when the network was
    down, some binaries and libraries are stored in `/bin`,
    `/sbin` and `/lib`
  \item `/bin` and `/sbin` contain programs like `ls`,
    `ip`, `cp`, `bash`, etc.
  \item `/lib` contains the C library and sometimes a few other
    basic libraries
  \item All other programs and libraries are in `/usr`
  \item Update: distributions are now making `/bin` link to
    `/usr/bin`, `/lib` to `/usr/lib` and `/sbin`
    to `/usr/sbin`. Details on
    \url{https://systemd.io/THE_CASE_FOR_THE_USR_MERGE/}.
  \end{itemize}

